% Options for packages loaded elsewhere
% Options for packages loaded elsewhere
\PassOptionsToPackage{unicode}{hyperref}
\PassOptionsToPackage{hyphens}{url}
%
\documentclass[
  ignorenonframetext,
  aspectratio=169,
  spanish,
]{beamer}
\newif\ifbibliography
\usepackage{pgfpages}
\setbeamertemplate{caption}[numbered]
\setbeamertemplate{caption label separator}{: }
\setbeamercolor{caption name}{fg=normal text.fg}
\beamertemplatenavigationsymbolshorizontal
% remove section numbering
\setbeamertemplate{part page}{
  \centering
  \begin{beamercolorbox}[sep=16pt,center]{part title}
    \usebeamerfont{part title}\insertpart\par
  \end{beamercolorbox}
}
\setbeamertemplate{section page}{
  \centering
  \begin{beamercolorbox}[sep=12pt,center]{section title}
    \usebeamerfont{section title}\insertsection\par
  \end{beamercolorbox}
}
\setbeamertemplate{subsection page}{
  \centering
  \begin{beamercolorbox}[sep=8pt,center]{subsection title}
    \usebeamerfont{subsection title}\insertsubsection\par
  \end{beamercolorbox}
}
% Prevent slide breaks in the middle of a paragraph
\widowpenalties 1 10000
\raggedbottom
\AtBeginPart{
  \frame{\partpage}
}
\AtBeginSection{
  \ifbibliography
  \else
    \frame{\sectionpage}
  \fi
}
\AtBeginSubsection{
  \frame{\subsectionpage}
}
\usepackage{iftex}
\ifPDFTeX
  \usepackage[T1]{fontenc}
  \usepackage[utf8]{inputenc}
  \usepackage{textcomp} % provide euro and other symbols
\else % if luatex or xetex
  \usepackage{unicode-math} % this also loads fontspec
  \defaultfontfeatures{Scale=MatchLowercase}
  \defaultfontfeatures[\rmfamily]{Ligatures=TeX,Scale=1}
\fi
\usepackage{lmodern}

\usetheme[]{Madrid}
\usecolortheme[]{seahorse}
\ifPDFTeX\else
  % xetex/luatex font selection
\fi
% Use upquote if available, for straight quotes in verbatim environments
\IfFileExists{upquote.sty}{\usepackage{upquote}}{}
\IfFileExists{microtype.sty}{% use microtype if available
  \usepackage[]{microtype}
  \UseMicrotypeSet[protrusion]{basicmath} % disable protrusion for tt fonts
}{}
\makeatletter
\@ifundefined{KOMAClassName}{% if non-KOMA class
  \IfFileExists{parskip.sty}{%
    \usepackage{parskip}
  }{% else
    \setlength{\parindent}{0pt}
    \setlength{\parskip}{6pt plus 2pt minus 1pt}}
}{% if KOMA class
  \KOMAoptions{parskip=half}}
\makeatother


\usepackage{longtable,booktabs,array}
\newcounter{none} % for unnumbered tables
\usepackage{calc} % for calculating minipage widths
\usepackage{caption}
% Make caption package work with longtable
\makeatletter
\def\fnum@table{\tablename~\thetable}
\makeatother
\usepackage{graphicx}
\makeatletter
\newsavebox\pandoc@box
\newcommand*\pandocbounded[1]{% scales image to fit in text height/width
  \sbox\pandoc@box{#1}%
  \Gscale@div\@tempa{\textheight}{\dimexpr\ht\pandoc@box+\dp\pandoc@box\relax}%
  \Gscale@div\@tempb{\linewidth}{\wd\pandoc@box}%
  \ifdim\@tempb\p@<\@tempa\p@\let\@tempa\@tempb\fi% select the smaller of both
  \ifdim\@tempa\p@<\p@\scalebox{\@tempa}{\usebox\pandoc@box}%
  \else\usebox{\pandoc@box}%
  \fi%
}
% Set default figure placement to htbp
\def\fps@figure{htbp}
\makeatother



\ifLuaTeX
\usepackage[bidi=basic,shorthands=off]{babel}
\else
\usepackage[bidi=default,shorthands=off]{babel}
\fi
\ifLuaTeX
  \usepackage{selnolig} % disable illegal ligatures
\fi


\setlength{\emergencystretch}{3em} % prevent overfull lines

\providecommand{\tightlist}{%
  \setlength{\itemsep}{0pt}\setlength{\parskip}{0pt}}



 


\makeatletter
\@ifpackageloaded{caption}{}{\usepackage{caption}}
\AtBeginDocument{%
\ifdefined\contentsname
  \renewcommand*\contentsname{Tabla de contenidos}
\else
  \newcommand\contentsname{Tabla de contenidos}
\fi
\ifdefined\listfigurename
  \renewcommand*\listfigurename{Listado de Figuras}
\else
  \newcommand\listfigurename{Listado de Figuras}
\fi
\ifdefined\listtablename
  \renewcommand*\listtablename{Listado de Tablas}
\else
  \newcommand\listtablename{Listado de Tablas}
\fi
\ifdefined\figurename
  \renewcommand*\figurename{Figura}
\else
  \newcommand\figurename{Figura}
\fi
\ifdefined\tablename
  \renewcommand*\tablename{Tabla}
\else
  \newcommand\tablename{Tabla}
\fi
}
\@ifpackageloaded{float}{}{\usepackage{float}}
\floatstyle{ruled}
\@ifundefined{c@chapter}{\newfloat{codelisting}{h}{lop}}{\newfloat{codelisting}{h}{lop}[chapter]}
\floatname{codelisting}{Listado}
\newcommand*\listoflistings{\listof{codelisting}{Listado de Listados}}
\makeatother
\makeatletter
\makeatother
\makeatletter
\@ifpackageloaded{caption}{}{\usepackage{caption}}
\@ifpackageloaded{subcaption}{}{\usepackage{subcaption}}
\makeatother

\usepackage{bookmark}
\IfFileExists{xurl.sty}{\usepackage{xurl}}{} % add URL line breaks if available
\urlstyle{same}
\hypersetup{
  pdftitle={El Barranco de Santos},
  pdfauthor={Análisis del Plan General de Ordenación},
  pdflang={es},
  hidelinks,
  pdfcreator={LaTeX via pandoc}}


\title{\texorpdfstring{El Barranco de Santos}{El Barranco de Santos}}
\subtitle{\texorpdfstring{De Cicatriz Urbana a Sistema General
Verde}{De Cicatriz Urbana a Sistema General Verde}}
\author{Análisis del Plan General de Ordenación}
\date{2026-05-08}

\begin{document}
\frame{\titlepage}


\begin{frame}{Introducción Histórica}
\protect\phantomsection\label{introducciuxf3n-histuxf3rica}
\begin{block}{El Barranco como Límite}
\protect\phantomsection\label{el-barranco-como-luxedmite}
\begin{itemize}
\tightlist
\item
  \textbf{La Cicatriz Física:} Históricamente, el Barranco de Santos ha
  funcionado como una barrera natural que divide Santa Cruz de Tenerife.
\item
  \textbf{Uso Tradicional:} Espacio marginal, infrautilizado y, en
  tramos, degradado.
\item
  \textbf{El Cambio de Paradigma:} El nuevo planeamiento busca
  transformar este ``vacío'' en el eje articulador de la ciudad.
\end{itemize}
\end{block}
\end{frame}

\begin{frame}{Gestión Urbanística}
\protect\phantomsection\label{gestiuxf3n-urbanuxedstica}
En el planeamiento español (y muy específicamente en el de Santa Cruz de
Tenerife), definen los \textbf{Sistemas de Actuación}. Básicamente,
indican qué mecanismo legal se usará para urbanizar ese suelo, quién
pone el dinero y quién lidera el proceso.

\begin{itemize}
\item
  \textbf{ASEXPR: Sistema de Expropiación} Es el sistema de iniciativa
  \textbf{pública} pura.
\item
  \textbf{ASCOOP: Sistema de Cooperación} Iniciativa \textbf{pública},
  pero financiada por los \textbf{privados}.
\end{itemize}

\begin{itemize}
\tightlist
\item
  \textbf{AACONV: Sistema de Concierto} Es un sistema \textbf{híbrido}
  (público-privado).
\end{itemize}

\begin{itemize}
\tightlist
\item
  \textbf{ASPRIV: Sistema Privado (generalmente Compensación)}
  Iniciativa \textbf{privada} bajo supervisión pública.
\end{itemize}
\end{frame}

\begin{frame}{El Futuro del Barranco}
\protect\phantomsection\label{el-futuro-del-barranco}
\begin{block}{Un Eje de Resiliencia}
\protect\phantomsection\label{un-eje-de-resiliencia}
\begin{itemize}
\tightlist
\item
  \textbf{Conectividad:} Un corredor verde que une las colinas con el
  mar.
\item
  \textbf{Sostenibilidad:} Protección de la flora autóctona y gestión de
  riesgos hidráulicos.
\item
  \textbf{Conclusión:} El Barranco de Santos deja de ser una división
  para ser el corazón verde de Santa Cruz.
\end{itemize}
\end{block}
\end{frame}

\begin{frame}{Gestion del Puente Zurita}
\protect\phantomsection\label{gestion-del-puente-zurita}
\begin{figure}[H]

{\centering \pandocbounded{\includegraphics[keepaspectratio]{PuenteZurita_GES2.png
}}

}

\caption{label}

\end{figure}%
\end{frame}

\begin{frame}{Gestion del Puente Serrador}
\protect\phantomsection\label{gestion-del-puente-serrador}
\begin{figure}[H]

{\centering \pandocbounded{\includegraphics[keepaspectratio]{PuenteSerrador_GES2.png
}}

}

\caption{label}

\end{figure}%
\end{frame}

\begin{frame}{Gestion del Puente Galceran}
\protect\phantomsection\label{gestion-del-puente-galceran}
\begin{figure}[H]

{\centering \pandocbounded{\includegraphics[keepaspectratio]{PuenteGalceran_GES2.png}}

}

\caption{Alt text}

\end{figure}%
\end{frame}

\begin{frame}{Puente Zurita Uso de Zonas}
\protect\phantomsection\label{puente-zurita-uso-de-zonas}
\begin{figure}[H]

{\centering \pandocbounded{\includegraphics[keepaspectratio]{PuenteZurita_ZUSO.png}}

}

\caption{Puente Zurita Uso de Zonas}

\end{figure}%
\end{frame}

\begin{frame}{Puente Galceran Uso de Zonas}
\protect\phantomsection\label{puente-galceran-uso-de-zonas}
\begin{figure}[H]

{\centering \pandocbounded{\includegraphics[keepaspectratio]{PuenteGalceran_ZUSO.png}}

}

\caption{Puente Galceran Uso de Zonas}

\end{figure}%
\end{frame}

\begin{frame}{{[}Puente Serrador Uso de Zonas}
\protect\phantomsection\label{puente-serrador-uso-de-zonas}
\begin{figure}[H]

{\centering \pandocbounded{\includegraphics[keepaspectratio]{PuenteSerrador_ZUSO.png}}

}

\caption{Puente Serrador Uso de Zonas}

\end{figure}%
\end{frame}

\begin{frame}{Estructuras}
\protect\phantomsection\label{estructuras}
\pandocbounded{\includegraphics[keepaspectratio]{Barranco de Santos Estructuras.png}}
\end{frame}

\begin{frame}{Fotos historicas}
\protect\phantomsection\label{fotos-historicas}
\begin{figure}[H]

{\centering \pandocbounded{\includegraphics[keepaspectratio]{Photos-3-001/La_Concepcion.jpg}}

}

\caption{Alt text}

\end{figure}%
\end{frame}

\begin{frame}{Fotos historicas}
\protect\phantomsection\label{fotos-historicas-1}
\begin{figure}[H]

{\centering \pandocbounded{\includegraphics[keepaspectratio]{Photos-3-001/LAVANDERAS DEL BARRANCO SANTOS.jpg
}}

}

\caption{label}

\end{figure}%
\end{frame}

\begin{frame}{Fotos historicas}
\protect\phantomsection\label{fotos-historicas-2}
\begin{figure}[H]

{\centering \pandocbounded{\includegraphics[keepaspectratio]{Photos-3-001/13006893.jpg
}}

}

\caption{label}

\end{figure}%
\end{frame}

\begin{frame}{Fotos historicas}
\protect\phantomsection\label{fotos-historicas-3}
\begin{figure}[H]

{\centering \pandocbounded{\includegraphics[keepaspectratio]{Photos-3-001/baranco santos y el hospital de san jose.jpg
}}

}

\caption{label}

\end{figure}%
\end{frame}

\begin{frame}{Fotos historicas}
\protect\phantomsection\label{fotos-historicas-4}
\begin{figure}[H]

{\centering \pandocbounded{\includegraphics[keepaspectratio]{Photos-3-001/Barranco de Santos 2.jpg
}}

}

\caption{label}

\end{figure}%
\end{frame}

\begin{frame}{Fotos historicas}
\protect\phantomsection\label{fotos-historicas-5}
\begin{figure}[H]

{\centering \pandocbounded{\includegraphics[keepaspectratio]{Photos-3-001/Barranco Santos 6.jpg
}}

}

\caption{label}

\end{figure}%
\end{frame}

\begin{frame}{Fotos historicas}
\protect\phantomsection\label{fotos-historicas-6}
\begin{figure}[H]

{\centering \pandocbounded{\includegraphics[keepaspectratio]{Photos-3-001/galderan sobre barranco santos.jpg}}

}

\caption{Alt text}

\end{figure}%
\end{frame}

\begin{frame}{Fotos historicas}
\protect\phantomsection\label{fotos-historicas-7}
\begin{figure}[H]

{\centering \pandocbounded{\includegraphics[keepaspectratio]{Photos-3-001/baranco_de_santos_foto.jpg}}

}

\caption{Alt text}

\end{figure}%
\end{frame}

\begin{frame}{Apéndice: Notas sobre los Datos}
\protect\phantomsection\label{apuxe9ndice-notas-sobre-los-datos}
\begin{itemize}
\tightlist
\item
  \textbf{Fuente:} Plan General de Ordenación (PGO).
\item
  \textbf{Herramientas:} Extracción mediante scripts de Python y
  visualización en QGIS.
\item
  \textbf{Gemini AI}
\end{itemize}
\end{frame}




\end{document}
